\section{Progetto}

Lo scopo del progetto è la costruzione di un programma \verb|C++| che simuli l'interazione tra due specie animali secondo il \href{https://it.wikipedia.org/wiki/Equazioni_di_Lotka-Volterra}{Modello Lotka-Volterra}.
Le due specie sono identificate tramite due variabili: $x$ e $y$, rispettivamente per le prede e i predatori. Gli ultimi si cibano dei primi.
\subsection{Descrizione del sistema}
Identificate le popolazioni di prede e predatori ad un certo istante temporale come $x(t)$ e $y(t)$, ne modelliamo l'andamento temporale tramite il set di equazioni differenziali, dette Equazioni di Lotka-Volterra:
\[
    \frac{\mathrm{d}x}{\mathrm{dt}} = \left[ A - B y(t)\right] x(t) \qquad \qquad    \frac{\mathrm{d}y}{\mathrm{dt}} = \left[ C x(t) - D \right] y(t)
\]
dove i coefficienti rappresentano
\begin{itemize}
    \item A - tasso di riproduzione delle prede
    \item B - tasso di mortalità delle prede
    \item C - tasso di riproduzione dei predatori
    \item D - tasso di mortalità dei predatori
\end{itemize}
La riproduzione è data dal consumo delle risorse disponibili alle due specie.
La mortalità delle prede è data dall'interazione con i predatori, che ne causano la morte; analogamente la riproduzione dei predatori è data dalla disponibilità di prede nel sistema e la mortalità dalla scarsità di prede.
Tutti questi parametri devono essere considerati maggiori di zero.

Durante la simulazione viene usata la versione discretizzata delle equazioni, tramite il \href{https://it.wikipedia.org/wiki/Metodo_di_Eulero_semi-implicito}{\textit{metodo di Eulero simplettico}}:

\[
    x_i = x_{i-1} + \left(A - B y_i\right) x_{i-1} \Delta t
    \qquad \qquad
    y_i = y_{i-1} + \left(C x_{i-1} - D\right) y_{i-1} \Delta t
\]
La scelta di moltiplicare $B$ per $y_i$ nel calcolo di $x_i$ si discosta dal metodo di integrazione che invece prevede $y_{i-1}$.
Questa scelta garantisce una stabilità maggiore del calcolo dell'integrale primo del moto $H(x,y)$, ma impone l'accorgimento di dover calcolare ad ogni passo dell'integrazione numerica prima $y_i$ e poi $x_i$.

\subsection{Dinamica del sistema}
Consideriamo uno stato iniziale determinato da $(x_0, y_0)$, tali che siano entrambi maggiori di 0, allora le orbite del sistame sono ristrette a valori positivi di $x(t)$ e $y(t)$. Questa richiesta è giustificata dal senso fisico del sistema: un numero di individui negativo non ha senso di esistere.

\par Risolvendo il sistema di equazioni differenziali si ottengo due punti di equilibrio:
\[
    e_1 = (0,0) \quad \mathrm{e} \quad e_2 = \left(\frac{D}{C}, \frac{A}{B} \right)
\]
che sono, rispettivamente, il caso con assenza contemporanea di prede e predatori e il secondo è l'unico caso non banale in cui l'evoluzione del sistema garantisce la stazionarietà delle popolazioni: si tratta di un perfetto equilibrio tra tutti i parametri del sistema.

Inoltre, si può ottenere una funzione che assume il ruolo di "energia" del sistema, per cui scelte delle condizioni si conserva lungo questa traiettoria. Per questa similitudine, verrà indicata con $H$:
\[
    H(x,y) = C x + B y - \left[D \ln(x) + A \ln(y)\right]
\]
In particolare il punto di equilibrio $\left(\frac{D}{C}, \frac{A}{B} \right)$ rappresenta anche il punto di minimo di $H(x,y)$

\subsection{Obiettivi del codice}
Forniti in input i valori iniziali $(x_0,y_0)$ e i quattro parametri delle equazioni di Lotka-Volterra, il programma sviluppa una simulazione che ad ogni step temporale calcola la terna $(x_i, y_i, H_i)$ e la salvi su un file \verb|.txt|.
